\documentclass{article}
\usepackage{graphicx} % Required for inserting images
\usepackage{parskip} % Adds space between paragraphs

\title{Cyber Threat Intelligence Sharing over DC-Networks}
\author{Simon Meyer}
\date{\today}

\begin{document}

\maketitle

\section{Introduction}

In an increasingly interconnected digital landscape, proactively defending against cyber threats presents a collective challenge that goes beyond the capabilities of individual organizations.
The exchange of Cyber Threat Intelligence (CTI) has proven essential for significantly increasing the detection rate of attacks.
Nevertheless, many stakeholders, particularly in the critical infrastructure sector, are reluctant to share sensitive threat information.
One reason for this is the concern about unintentionally revealing their own identity or infrastructure details, which could lead to reputational damage or targeted follow-up attacks.

Existing anonymous CTI-sharing systems often protect the anonymity of the CTI provider only from the system operators.
In some cases, anonymity is also maintained vis-à-vis other participants, but only if they are honest and follow the protocol.
These systems are typically based on blockchain technology to preserve anonymity.
However, these approaches do not protect against network observers, as a participant in the system who monitors the network interfaces can link the transmission of data from one participant to the receipt of CTI.

The scientific literature is increasingly addressing this trade-off between information exchange and data protection through privacy-enhancing technologies (PETs).
While conventional anonymization services such as mix networks and onion routing approaches like Tor form a widely used foundation for anonymous communication, their architecture makes them unsuitable for CTI sharing.

As a robust alternative, Dining Cryptographers Networks (DC-nets), based on the theoretical foundations of David Chaum, offer information-theoretically guaranteed sender anonymity.
Unlike mix networks, protection in DC networks does not rely on computational complexity or packet delay, but on indistinguishability within a logical group through cryptographic overlay.

This paper presents a system design for the anonymous exchange of CTI data that leverages the strong anonymity of DC-nets.
The goal is to create a framework that enables organizations to share threat intelligence in real time without compromising the confidentiality of the sender’s identity.


\end{document}

